\documentclass[10pt,a4paper]{amsart}
\usepackage[margin=19mm]{geometry}
\usepackage{amsmath,amssymb,mathtools}
\usepackage[scr=rsfs,cal=euler]{mathalfa}
\usepackage{xfrac}
\usepackage[unicode,hidelinks,bookmarks=false]{hyperref}
\newcommand{\ie}{i.e.\ }
\newcommand{\cf}{cf.\ }
\newcommand{\resp}{respectively\ }
\newcommand{\Subsection}{\subsection}
\providecommand{\nobreakdash}{\nobreak\hbox{-}}
\setlength{\emergencystretch}{2em}
\multlinegap0pt
\usepackage{amssymb}
\usepackage{mathtools}
\newtheorem{theorem}{Theorem}
\newtheorem{lemma}{Lemma}
\newcommand{\CH}{\mathbb{CH}}
\newcommand{\Hess}{\operatorname{Hess}}
\newcommand{\II}{\mathrm{II}}
\newcommand{\PU}{\mathrm{PU}}
\newcommand{\Teich}{\mathcal T}
\newcommand{\dd}{\,\mathrm d}
\newcommand{\im}{\operatorname{im}}
\newcommand{\tr}{\operatorname{tr}}
\title{Critical-point-free energy for\\ fractional-Toledo representations}
\author{Rivu Bardhan, Anu Dhochak, Pradip Kumar}
\date{Source: arXiv:2608.15714v1 (CC BY 4.0)}
\begin{document}
\maketitle

\noindent\emph{Source and licence.} Critical-point-free energy for\\ fractional-Toledo representations. Version arXiv:2608.15714v1 (CC BY 4.0), distributed by arXiv under the Creative Commons Attribution 4.0 International licence, http://creativecommons.org/licenses/by/4.0/, as recorded in the arXiv licence metadata for this exact version and shown on the abstract page https://arxiv.org/abs/2608.15714v1. Full licence notice: \texttt{licenses/arxiv-2608.15714-CC-BY-4.0.txt}.\par\smallskip

\noindent\textit{Editorial scope.} Counted unit: the article's theorem thm:main (source0.tex lines 108--122). The article's complete original proof of it is delivered with it (source0.tex lines 673--754, one \texttt{proof} environment of 82 lines, which the article places away from the statement). No mathematics is written, repaired or re-derived: the statement and the proof are the article's own lines, in the article's own notation, and no retained line is substituted or re-flowed.  Retained uncounted as the source-local context the proof needs: lines 181--186; lines 533--548; lines 614--623.  Nothing was omitted from any retained span, so the package carries no omission record.  No retained line contains a citation, so the wrapper carries no bibliography and no bibliography entry.  No retained line contains a figure environment, an image-inclusion command or drawing code, so no image is delivered and the package declares no asset dependency.  Editorial formatting only, in addition: an AMS \texttt{amsart} preamble that restates the article's own declarations and macros used by the retained lines, namely the environments \texttt{theorem}, \texttt{lemma} on separate counters, together with the standard packages those lines need.  The retained environments print in this document as Theorem 1, Lemma 1 in order of appearance, whereas the article prints the same items with the numbering of its own full document.  The complete native source of the article as distributed in the arXiv e-print for this exact version is bundled unchanged as \texttt{sources/207774/source0.tex}.
\begin{equation}\label{eq:critical-equivalence}
 J\in\operatorname{Crit}(E_\rho)
 \quad\Longleftrightarrow\quad q_{\rho,J}=0
 \quad\Longleftrightarrow\quad f_{\rho,J}
 \text{ is weakly conformal}.
\end{equation}

\begin{lemma}\label{lem:tube}
Let $X$ be a Hadamard manifold with $K_X\le-1$, and let
$Y^2\subset X$ be a complete embedded minimal surface such that
\[
 \sup_{y\in Y}\|\II_Y\|_{\mathrm{op}}\le\eta<\tfrac12,
 \qquad
 \exp^\perp:NY\longrightarrow X\ \text{is a diffeomorphism}.
\]
Put $\Psi=d(\,\cdot\,,Y)^2$.  Then $\Psi\in C^\infty(X)$ and, for every
$R>0$, there is $c_R>0$ such that, whenever $d(x,Y)\le R$ and
$L\subset T_xX$ is a two-plane,
\begin{equation}\label{eq:tube-estimate}
 \tr_L\Hess\Psi\ge c_R\Psi(x).
\end{equation}
Here $\tr_L$ denotes the trace of the restriction to $L$.
\end{lemma}

\begin{lemma}[Branched trapping]\label{lem:trapping}
Assume that $X$ and $Y$ satisfy the hypotheses of Lemma~\ref{lem:tube}.  Let $M$
be a closed connected Riemann surface, let
$\sigma:\pi_1(M)\to\operatorname{Isom}(X)$ have image preserving $Y$, and let
\[
 f:\widetilde M\longrightarrow X
\]
be a nonconstant $\sigma$-equivariant harmonic map which is weakly conformal.
Then $f(\widetilde M)\subset Y$.
\end{lemma}

\begin{theorem}\label{thm:main}
Fix a positive integer $d$ with $3\nmid d$.  There exists $h_0(d)$ such that,
for every $h\ge h_0(d)$ and every $g>h$, there is an irreducible reductive
representation
\[
 \rho_{g,h,d}:\pi_1(S_g)\longrightarrow\PU(2,1)
\]
satisfying
\begin{equation}\label{eq:toledo-main}
 \tau(\rho_{g,h,d})=2h-2-\frac{2d}{3}\notin\mathbb Z,
 \qquad
 \operatorname{Crit}(E_{\rho_{g,h,d}})=\varnothing.
\end{equation}
Branch points are allowed in this assertion.
\end{theorem}

\begin{proof}[Proof of Theorem~\ref{thm:main}]
Suppose, for contradiction, that $J\in\Teich(S_g)$ is a critical point of
$E_\rho$.  By the equivalences in~\eqref{eq:critical-equivalence}, the
Corlette harmonic map
\[
f=f_{\rho,J}:(\widetilde S_g,J)\longrightarrow\CH^2
\]
is weakly conformal and is therefore a possibly branched minimal map.  It is
nonconstant: otherwise $\im\rho$ would fix its value, contrary to
irreducibility.

The group $\im\rho=\im\rho_0$ preserves the almost-Fuchsian plane
$Y=F(\widetilde S_h)$.  Applying Lemma~\ref{lem:trapping} yields
\begin{equation}\label{eq:image-in-Y}
 f(\widetilde S_g)\subset Y.
\end{equation}
Since $F:\widetilde S_h\to Y$ is an embedding, define
\[
 H=F^{-1}\circ f:\widetilde S_g\longrightarrow\widetilde S_h.
\]
Indeed, injectivity of $F$ and equivariance give
\[
 F(H(\gamma x))=f(\gamma x)=\rho_0(P_*(\gamma))f(x)
 =F(P_*(\gamma)H(x)).
\]
Therefore
\begin{equation}\label{eq:H-equivariance}
 H(\gamma x)=P_*(\gamma)H(x),
 \qquad \gamma\in\pi_1(S_g).
\end{equation}
Consequently $H$ descends to a smooth map
\begin{equation}\label{eq:Hbar}
 \overline H:(S_g,J)\longrightarrow(S_h,J_0)
\end{equation}
which induces $P_*$ on fundamental groups after compatible basepoint choices
(and induces it up to an inner automorphism without such choices).

Equip $\widetilde S_h$ with the induced metric $g_Y=F^*g_{\CH^2}$.  For each
$\gamma\in\pi_1(S_h)$, equivariance and invariance of the ambient metric give
\[
 \gamma^*g_Y=(F\circ\gamma)^*g_{\CH^2}
 =(\rho_0(\gamma)\circ F)^*g_{\CH^2}=g_Y.
\]
Thus $g_Y$ descends to a metric on $S_h$.  Because $F$ is holomorphic, the
conformal class of $g_Y$ is $J_0$.  Since $F$ is an isometric immersion for
$g_Y$ and $f=F\circ H$ is weakly conformal, $H$ is weakly conformal.

For a local orthonormal frame $e_1,e_2$ on the domain, the composition
formula for tension fields is
\begin{equation}\label{eq:tension-composition}
 \tau(F\circ H)=\dd F(\tau(H))+
 \sum_{i=1}^2\II_F\bigl(\dd H(e_i),\dd H(e_i)\bigr).
\end{equation}
Where $\dd H\ne0$, choose an orthonormal frame with
$\dd H(e_i)=\lambda E_i$ and $E_1,E_2$ orthonormal.  Then the last term is
$\lambda^2(\II_F(E_1,E_1)+\II_F(E_2,E_2))$ and vanishes because $F$ is
minimal.  Where $\dd H=0$, it vanishes directly.  Since $F\circ H=f$ is
harmonic and $\dd F$ is injective, equation~\eqref{eq:tension-composition}
implies
\[
 \tau(H)=0.
\]
Thus $\overline H$ is a nonconstant harmonic weakly conformal map between
Riemann surfaces.  Its branch set is discrete, its complement is connected,
and the orientation sign is constant there.  The standard local
characterization therefore shows that $\overline H$ is globally holomorphic
or globally antiholomorphic.

The closed surfaces $S_g$ and $S_h$ are aspherical.  Because
$\overline H_*=P_*$, the maps $\overline H$ and $P$ are homotopic and
\[
 \deg\overline H=\deg P=1.
\]
An antiholomorphic nonconstant map has negative signed degree, so
$\overline H$ is holomorphic.  A nonconstant holomorphic map between compact
connected Riemann surfaces is surjective, and the local degrees over any
fiber are positive integers whose sum is the global degree.  Since
$\deg\overline H=1$, every fiber consists of one point of local degree one.
Thus $\overline H$ is unramified and bijective, hence biholomorphic.
Consequently $g=h$, contradicting $g>h$.  Therefore $E_\rho$ has no critical
point.
\end{proof}

\end{document}
